\subsection{AFFINE WEYL GROUP}

For \(\alpha \in \mathbb{R}\) and \(k\in \mathbf{Z}\) , let \(\mathrm{L}_{\alpha ,k}\) be the hyperplane of E defined by:

\[
\mathrm{L} _ {\alpha , k} = \{x \in \mathrm{E} | \langle \alpha , x \rangle = k \}
\]

and let \(s_{\alpha, k}\) be the orthogonal reflection with respect to \(L_{\alpha, k}\) . Then

\[
s _ {\alpha , k} (x) = x - (\langle \alpha , x \rangle - k) \alpha^ {\prime} = s _ {\alpha , 0} (x) + k \alpha^ {\prime}
\]

for all \(x \in E\) . In other words,

\[
s _ {\alpha , k} = t (k \alpha^ {\prime}) \circ s _ {\alpha}\tag{1}
\]

where \(s_{\alpha}\) is the orthogonal reflection with respect to the hyperplane \(L_{\alpha} = L_{\alpha,0}\) , i.e.~the reflection associated to the root \(\alpha\) .

Formula (1) shows that \(s_{\alpha, k}\) does not depend on the choice of scalar product.

DEFINITION 1. The group of affine transformations of E generated by the reflections \(s_{\alpha,k}\) for \(\alpha \in R\) and \(k \in Z\) is called the affine Weyl group of the root system R and denoted by \(\mathrm{W}_{a}(\mathrm{R})\) (or simply by \(W_{a}\) ).

PROPOSITION 1. The group \(\mathrm{W}_a\) is the semi-direct product of \(\mathrm{W}\) by \(\mathrm{Q}\) .

Since W is generated by the reflections \(s_{\alpha}\) , it is contained in \(W_{a}\) . On the other hand, \(t(\alpha^{\prime}) = s_{\alpha,1} \circ s_{\alpha}\) if \(\alpha \in R\) , which shows that \(Q \subseteq W_{a}\) .

Since W leaves \(Q(R^{\prime})\) stable ( \(\S\) 1, no. 9), the group G of affine transformations generated by W and Q is the semi-direct product of W by Q. Now \(G \subseteq W_{a}\) from above and \(s_{\alpha,k} \in G\) for all \(\alpha \in R\) and \(k \in Z\) by (1). It follows that \(W_{a} = G\) .

PROPOSITION 2. The group \(W_{a}\) , with the discrete topology, acts properly on E and permutes the hyperplanes \(L_{\alpha,k}\) (for \(\alpha \in R\) and \(k \in Z\) ).

Since \(Q(R^{\prime})\) is a discrete subgroup of \(V^{*}\) , the group Q acts properly on E. Hence, so does \(W_{a} = W.Q\) , since W is finite. Moreover, for \(\alpha, \beta \in R\) and \(k \in Z\) :

\[
s _ {\beta} (\mathrm{L} _ {\alpha , k}) = \mathrm{L} _ {\gamma , k} \text {with} \gamma = s _ {\beta} (\alpha) \in \mathrm{R},
\]

\[
t (\beta^ {\prime}) \left(\mathrm{L} _ {\alpha , k}\right) = \mathrm{L} _ {\alpha , k + n (\alpha , \beta)},
\]

where \(n(\alpha, \beta) = \langle \beta^{\prime}, \alpha \rangle\) is an integer, hence the second assertion.

We can thus apply the results of Chapter V, § 3 to \(W_{a}\) acting on E. To avoid any confusion with the chambers of the Weyl group W in \(V^{*}\) , we shall call the chambers determined by the system of hyperplanes \(L_{\alpha,k}\) (for \(\alpha \in R\) and \(k \in Z\) ) in E alcoves. The group \(W_{a}\) thus acts simply transitively on the set of alcoves and the closure of an alcove is a fundamental domain for \(W_{a}\) acting on E (Chap. V, § 3, no. 2, Th. 1 and no. 3, Th. 2). It is clear that the Weyl group W is identified with the canonical image \(U(W_{a})\) of \(W_{a}\) in the orthogonal group of \(V^{*}\) (cf.~Chap. V, § 3, no. 6). It follows that \(W_{a}\) is essential (Chap. V, § 3, no. 7) and that \(W_{a}\) is irreducible if and only if the root system R is (§ 1, no. 2, Cor. of Prop. 5). If R is irreducible, every alcove is an open simplex (Chap. V, § 3, no. 9, Prop. 8). In the general case, the canonical product decomposition of the affine space E (Chap. V, § 3, no. 8) corresponds to the decomposition of R into irreducible components. In particular, the alcoves are products of open simplexes.

Note also that the Cor. of Th. 1 of Chap. V, § 3, no. 2 shows that the \(s_{\alpha,k}\) are the only reflections in \(\mathbf{W}_a\) .

